\section{Picture dictionaries and multimodal practice}

\subsection{Why picture dictionaries matter}

Picture dictionaries are a useful point of contact between C-LARA-2's
multimodal authoring infrastructure and the needs of minority and Indigenous
language communities. Picture-based resources can support learners with
limited literacy, communities whose orthographies are still being consolidated,
and contributors who want to create accessible resources for younger learners.
They also make a demanding requirement especially clear: community members
must be able to author and review individual lexical entries without having to
master the technical workflow that underlies a general-purpose annotated text.

The original technical motivation for adding picture dictionaries was picture
glossing. In this learner-facing functionality, a word or expression in a
C-LARA-2 text is associated with a dictionary entry; clicking it can display
the corresponding picture alongside other support such as audio playback and
concordance information. The required infrastructure is a reviewed, reusable
association between lexical items and images. Picture glossing has now been
implemented, but has not yet been tested in learner use.

In practice, the first substantial demand for picture dictionaries has come
from collaborators working in Indigenous community contexts, where the
immediate use cases are not identical to the original picture-glossing
scenario. Contributors need accessible ways to create, inspect and control
picture-based lexical resources and to reuse them in exercises. These needs
made the dictionary-entry workflow, including visible cultural guidance and
human review of images, the more urgent design problem.

A picture dictionary is represented internally as an ordinary C-LARA-2 project
whose text has one lexical item or lexicalised expression per page. This was a
useful reuse decision: it made existing project, annotation, image and
compilation infrastructure available without introducing a separate data
model. It did not, however, imply that the standard project interface was the
right interface for dictionary construction.

\subsection{Why the general-purpose workflow was not enough}

The annotation interface described in Section~4 is designed for an extended
text containing pages, segments and token-level linguistic analysis. The image
workflow in Section~5 is designed for page-level scenes whose style and
recurring visual elements must be coordinated. Both workflows are appropriate
for their intended tasks. A picture dictionary is more constrained: the
relevant unit is a lexical entry, which needs a word, perhaps a lemma and part
of speech, a gloss or translation, optional human guidance, an image prompt
and a selected image considered together.

The initial picture-dictionary design exposed the underlying project as a
sequence of existing annotation and page-image operations. Contributors had to
move between separate views to manage information that conceptually belonged
to one entry. This created complexity without a corresponding benefit,
particularly where contributors were community members rather than technical
specialists, or where the project language could not itself be processed
reliably by the available AI model. In such cases, a human-provided English or
French translation may be needed to mediate prompt creation, while cultural
notes or restrictions must remain visible to the person reviewing the image.

\subsection{A dictionary-entry-centred interface}

Following discussion with intended users, C-LARA-2 retained the project-based
representation but replaced the user-experience layer with a specialised
picture-dictionary editor. Each lexical entry is displayed as an integrated
block containing the word, lemma, part of speech, gloss or translation,
optional suggestion, image-generation prompt and selected image. Contributors
can inspect and edit each component directly; they can also select entries in
bulk to create missing translations, prompts, images or other information.
The interface thus exposes the relevant parts of the annotation and image
workflows without asking a contributor to navigate the complete pipelines for
every word.

\begin{figure}[htbp]
\centering
\includegraphics[width=\linewidth]{images/picture_dictionary_editor.jpg}
\caption{Part of the specialised C-LARA-2 picture-dictionary editor. Each
lexical entry brings together word, lemma/POS information, gloss or
translation, optional human guidance, image-generation prompt and selected
image. Controls support review and batch generation for selected entries. A
French example is shown rather than community-language data because work
involving Indigenous languages may be subject to cultural-sensitivity and
data-sovereignty constraints.}
\label{fig:picture-dictionary-editor}
\end{figure}

The redesign was not merely cosmetic. The earlier interface treated dictionary
construction as a sequence of operations inherited from the general project
workflow; the replacement treats the dictionary entry as the organising unit.
Codex produced a usable implementation within a few hours once the revised
conceptual model had been agreed, and the consortium members involved in
testing responded positively. This is a useful example of AI-centred
development under human supervision: the AI did not independently identify the
community need or diagnose the usability problem; human collaborators supplied
those judgements, while Codex carried out the rapid operational redesign.

The episode also illustrates a practical consequence of AI-mediated
development. Conventional software work is subject to a sunk-cost pressure:
prior effort can make it difficult to abandon an unsuitable design
\citep{ArkesBlumer1985}. Here, the project was able to preserve the useful
underlying representation while discarding and rebuilding the user-facing
layer. The image-generation process itself still requires further work and
careful human review; rapid implementation is not evidence that generated
images are automatically pedagogically or culturally appropriate.

\subsection{Current exercises and future picture glossing}

At present, the picture-dictionary infrastructure is being used most
concretely to support picture-based exercises in Indigenous community
contexts. Picture glossing is also implemented, but should be treated as an
untested next application rather than as a demonstrated educational result.
The same community-authored lexical entries can support more than a dictionary
display. In picture glossing, clicking a word or expression in a C-LARA-2 text
can show its picture, play audio and provide concordance examples. The
platform can also create picture-to-word and word-to-picture multiple-choice
exercises, word scrambles and simple crosswords using picture clues. This
reuse matters: effort invested in reviewing a lexical entry can support
several forms of multimodal practice rather than one isolated artefact.

\begin{figure}[htbp]
\centering
\includegraphics[height=0.72\textheight]{images/word_scramble.jpg}
\caption{A picture-clue word-scramble exercise generated from a French picture
dictionary. The learner selects adjacent letters in the grid to spell the
lexical item represented by the image; two target words have already been
found.}
\label{fig:word-scramble}
\end{figure}

The word-scramble example is intentionally modest, but it makes the point
concrete. The learner receives a visual clue, performs an active lexical task,
receives feedback and moves to the next item. The relevant innovation is not
that a word scramble is novel in itself; it is that a reviewed dictionary
entry can be reused rapidly as one component of an integrated, multimodal
learning resource.

The project intends to return to the original picture-glossing idea in a
context involving multimodal texts for children or beginner learners. For
picture glossing to be useful throughout a text, a dictionary is likely to
need coverage of at least several hundred common words and expressions. Most
of the operational work can in principle be automated: candidate vocabulary
can be identified, translations and prompts proposed, candidate images
generated, and dictionary entries linked to textual tokens. Human review
remains necessary, but it can focus on accepting, correcting or rejecting
proposed entries rather than manually constructing the whole resource. This
is a promising area for pedagogical discussion with picture-based-learning
specialists, including Christ\`ele Maizonniaux, whose research interests are
closely relevant to the design choices involved.
